\begin{abstract}
\noindent\textbf{Question.} How does the number of threads affect the runtime, speedup and efficiency of a shared-memory dense \texttt{float64}
matrix multiplication on a 4-core / 8-thread laptop CPU, and from which size does parallel execution pay off?
\textbf{Method.} The blocked kernel of Assignment~2 (tile 96) is the sequential baseline; the parallel version splits the rows of $C$ among
OpenMP threads and runs the same kernel, so both perform bit-identical arithmetic on the same inputs. Thread counts 1--8 plus 16 (oversubscription),
$n=100$--$4000$, monotonic kernel timer, warm-up, 5 samples (3 at $n=4000$), median and IQR, randomised order, a repeated pass, \texttt{ru\_maxrss}
memory, budgets of 120~s and 4~GiB, 1404 correctness checks and per-run validation; an optional extension compares scalar, SSE2 and AVX2 builds.
\textbf{Results.} All 110 runs were correct and bit-identical to the baseline. Speedup is almost linear up to the 4 physical cores
($S(4)=3.40$--$3.95$, efficiency 0.85--0.99) and saturates there: 8 threads change it by $-15$ to $+6\%$ and 16 threads give no consistent gain.
Parallelism already paid off at $n=100$ ($S(4)=3.54$). With one thread the OpenMP version matches the baseline within the between-process noise, and
memory does not grow with the thread count. AVX2 is 3.9--5.1$\times$ faster than scalar code at equal thread count.
\textbf{Scope.} One laptop, unpinned threads, $\pm10\%$ variability between periods of the campaign and no hardware counters; conclusions apply to
this kernel, timing boundary and machine.
\end{abstract}
